% ============================================================
\section{Thiết lập thực nghiệm}
% ============================================================
\textbf{Bộ dữ liệu.} Thực nghiệm tiến hành trên \emph{VietnamNet}, gồm 2.000 bài
báo phân bố đều trên 20 chuyên mục (100 bài/chuyên mục), mỗi bài 500--800 từ kèm
tiêu đề, sapo (đoạn dẫn nhập) và danh sách keyphrase. Đánh giá tóm tắt, trích xuất keyphrase và phân loại đều thực hiện trên toàn bộ
2.000 bài. Tùy tác vụ, các thành phần được
dùng làm tham chiếu: tiêu đề + sapo làm tóm tắt tham chiếu; danh sách keyphrase do
LLM (DeepSeek~V4) sinh làm tham chiếu keyphrase~\cite{pavlovic2024effectiveness};
và 20 chuyên mục gốc làm
\emph{coarse-label proxy} cho không gian nhãn (do không tồn tại ground-truth
tuyệt đối~\cite{wan2024tnt}).

\textbf{Môi trường \& tham số.} Toàn bộ thực nghiệm chạy trên máy cá nhân CPU
Intel Core i7 thế hệ 12, 16\,GB RAM, không GPU; các mô hình embedding chạy trên
CPU qua \texttt{sentence-transformers}; LLM truy cập qua Google API. Tham số cố
định: Tầng~1 ($d=0{,}85$, $\varepsilon=10^{-5}$, 40\% số câu, tối thiểu 5); Tầng~2
(n-gram $1$--$3$, $m=10$, MMR $\lambda=0{,}5$); Tầng~3 (\texttt{gemini-2.5-flash},
Temperature $0{,}1$ cho phân loại và $0{,}3$ cho sinh nhãn, gửi tối đa 5
keyphrase/bài).

\textbf{Baseline.} Đối với tóm tắt: Lead-K, TextRank gốc (Word Overlap), TextRank
+ TF-IDF. Đối với keyphrase: RAKE, YAKE, TextRank từ khóa, TopicRank, Multipartite,
WordAttractionRank, KeyBERT + embedding đối xứng. Đối với khám phá không gian nhãn
end-to-end: \emph{X-MLClass-Gemini} -- bản cài đặt của X-MLClass~\cite{li2024open} giữ nguyên
mô tả gốc (chunk 50 token, 4 keyphrase/chunk, BERTopic, NLI entailment, 10 vòng
lặp long-tail). Chúng tôi gọi baseline này là X-MLClass-Gemini vì nó tuân theo
pipeline X-MLClass và chỉ thay LLM gốc bằng Gemini 2.5 Flash để đồng nhất backbone
khi so sánh.

\textbf{Độ đo.} Tóm tắt dùng ROUGE-1/2/L (F1). Keyphrase dùng Precision/Recall/F1
tại $k=5,10$ với cơ chế so khớp kết hợp \emph{exact match} và \emph{soft match}
(cosine $\geq 0{,}7$). Không gian nhãn được phân tích định tính bằng cách xếp mỗi
nhãn vào một trong năm nhóm thủ công -- \emph{Coarse}, \emph{Long-tail}, \emph{Entity},
\emph{Meta/Noise}, \emph{Cross-topic} -- từ đó trích ba chỉ số: tỉ lệ Coarse, tỉ
lệ Nhiễu (Entity+Meta/Noise) và số chuyên mục được phủ. Chất lượng gán nhãn đầu
cuối dùng Precision-at-$k$:
\begin{equation}
P@k = \frac{1}{N}\sum_{i=1}^{N}
\frac{\lvert C_i^{rk}\cap L_i\rvert}{\min(k,\lvert L_i\rvert)},
\label{eq:pak}
\end{equation}
trong đó $C_i^{rk}$ là $k$ nhãn dự đoán hàng đầu và $L_i$ là tập nhãn tham chiếu.
Do thiếu nhãn đa nhãn ground-truth, chuyên mục gốc được dùng làm $L_i$ với
$\lvert L_i\rvert=1$.

